\section{Derivations}
\label{app:derivations}
\subsection{Exact transition kernels of the SDE sampler}
ELF's stochastic sampler at step $k$ with step size $h_k$ and churn $\gamma$ predicts $\hat x_k=D_\theta(y_{k-1},t_{k-1},c)$, moves towards it, and re-noises: $y_k=(1-\gamma h_k)\,\hat x_k+\gamma h_k c_\sigma\,\varepsilon_k$ with $\varepsilon_k\sim\mathcal N(0,I)$ and $c_\sigma$ the noise scale of the interpolant at the target time (we absorb the deterministic drift terms into $D_\theta$; the released code is the reference). Indexing the chain by the post-noising iterates $y_k$ gives $p_\theta(y_k\mid y_{k-1})=\mathcal N(y_k;\alpha_kD_\theta(y_{k-1}),\sigma_k^2I)$ with $\alpha_k=1-\gamma h_k$ and $\sigma_k=\gamma h_kc_\sigma$, and the final decode $p(w\mid y_K)=\prod_i\mathrm{Cat}(w_i;\mathrm{softmax}(E\,D_\theta(y_K)_i))$ with $E$ the embedding matrix. The trajectory log-likelihood is the sum of these terms and requires no bound. We verified the implementation on CPU: recomputing the log-probability of a stored trajectory under the same parameters agrees with the value recorded at rollout to relative error $7\times10^{-5}$ (flow steps) and $10^{-3}$ (decode). Two caveats that matter in practice: $\gamma h_k\ge1$ makes $\alpha_k=0$ and every parameter gradient exactly zero (we assert $\alpha_k>0$), and storing states in bfloat16 perturbs summed log-ratios by $O(1)$ over $L\times d$ dimensions, so ratio-based estimators need per-dimension normalisation or fp32 storage.

\subsection{Score-function and pathwise gradients}
\label{app:variance}
For~(\ref{eq:chain}), $\nabla_\theta\log\mathcal N(y_k;\alpha_kD_\theta,\sigma_k^2I)=\sigma_k^{-2}(y_k-\alpha_kD_\theta)^{\top}\alpha_k\partial D_\theta/\partial\theta=\sigma_k^{-1}\varepsilon_k^{\top}\alpha_k\,\partial D_\theta/\partial\theta$, which is~(\ref{eq:sf}). The estimator's variance is that of a sum of $K\cdot L\cdot d$ terms each proportional to an independent standard normal, weighted by an advantage that depends on the trajectory only through the final reward. Write the reward as a function of the noise path, $r=r(\varepsilon_{1:K})$; the score-function estimate of $\partial_\theta\mathbb E[r]$ is $\mathbb E[r\sum_k\sigma_k^{-1}\varepsilon_k^{\top}\alpha_kJ_k]$ with $J_k=\partial D_\theta/\partial\theta$, so a coordinate contributes signal only through the covariance of $r$ with that coordinate's noise. In an AR model the analogous decomposition has one term per token and every term is a categorical decision the reward depends on; in the flow model, 86\% of positions are decided at the first step and never change (\S\ref{app:commit}), so almost every coordinate at almost every step has vanishing covariance with $r$ while contributing its full variance. With $N$ trajectories the signal scales as $N$ times the few informative covariances and the noise as $\sqrt{N\cdot K L d}$; the ratio of informative to total coordinates ($\sim$30 semantic decisions against $5\times10^{5}$--$2\times10^{6}$ coordinates) is the $10^{4}$--$10^{5}$ gap quoted in \S\ref{sec:estimators}. Refining the advantage (per-step, per-position, oracle) changes which coordinates receive weight but not this ratio, which is why Table~\ref{tab:sf_variants} is flat.

The pathwise estimator uses the same reparameterisation: $y_k$ is a deterministic function of $(\theta,\varepsilon_{1:k})$, so for a differentiable surrogate $V_\phi$ of $\mathbb E[r\mid y_k]$, $\partial_\theta V_\phi(y_k)=\partial_{y}V_\phi\cdot\alpha_k\,\partial_\theta D_\theta(y_{k-1})$ holds coordinate-wise with no noise factor. Its bias is that of the critic, which is why the critic must track the policy. We use one random $k$ per trajectory (the ``one-step'' variant) rather than back-propagating through the whole chain as in \citet{prabhudesai2023alignprop,clark2023draft}; with a critic on intermediate states the chain need not be differentiated, and memory is constant in $K$.

\subsection{Why the pathwise update ascends the reward: derivation of~(\ref{eq:pw})}
\label{app:pw_derivation}
This section derives the pathwise update of \S\ref{sec:estimators} from the objective~(\ref{eq:objective}) and makes precise in what sense the one-step \jpc{} update of \S\ref{sec:method} is a policy-gradient step on the expected reward. Fix a prompt $c$ and drop it from the notation. The chain~(\ref{eq:chain}) is $y_0=c_\sigma\varepsilon_0$ and $y_k=\alpha_kD_\theta(y_{k-1})+\sigma_k\varepsilon_k$ for $k=1,\dots,K$ with $\varepsilon_0,\dots,\varepsilon_K$ independent standard Gaussians, and the reward is a function of the last state, $r(y_K):= r(w(y_K),c)$, where $w(y_K)$ is the decoded text (the final ODE step and the decoder head are absorbed into this map).

\paragraph{Value functions.}
For $k=0,\dots,K$ define the value of a state at flow step $k$ under the current policy,
\begin{equation}
V_k(y):=\mathbb E\big[r(y_K)\,\big|\,y_k=y\big],
\label{eq:valuedef}
\end{equation}
where the expectation is over the remaining noise $\varepsilon_{k+1},\dots,\varepsilon_K$. Then $V_K=r$, the objective is $J(\theta)=\mathbb E[V_0(y_0)]$, and the tower property gives the one-step recursion
\begin{equation}
V_{k-1}(y)=\mathbb E_{\varepsilon_k}\big[V_k\big(\alpha_kD_\theta(y)+\sigma_k\varepsilon_k\big)\big].
\label{eq:bellman}
\end{equation}
Every $V_k$ with $k<K$ depends on $\theta$ through the downstream transitions; $V_K=r$ does not. The 1D critic of \S\ref{sec:method} is a learned approximation of $V_k$, and the 2D critic is a learned approximation of $V_k$ that is additive over the positions of $y$.

\paragraph{Step 1: the policy gradient is a sum over transitions.}
$J$ depends on $\theta$ through each of the $K$ transitions. Give the $k$-th transition its own copy $\theta_k$ of the parameters, so that $J=J(\theta_1,\dots,\theta_K)$ with all copies tied to $\theta$; by the chain rule,
\begin{equation}
\nabla_\theta J=\sum_{k=1}^{K}\frac{\partial J}{\partial\theta_k}\Big|_{\theta_1=\dots=\theta_K=\theta}.
\label{eq:sumover}
\end{equation}
This is the decomposition behind every per-step estimator, including the score-function form~(\ref{eq:sf}): writing the $k$-th transition density as $\mathcal N(y_k;\alpha_kD_{\theta_k}(y_{k-1}),\sigma_k^2I)$ and applying the likelihood-ratio identity to the expectation of the downstream value gives
\begin{equation}
\frac{\partial J}{\partial\theta_k}=\mathbb E\Big[V_k(y_k)\,\nabla_{\theta_k}\log\mathcal N\big(y_k;\alpha_kD_{\theta_k}(y_{k-1}),\sigma_k^2I\big)\Big]=\mathbb E\Big[V_k(y_k)\,\frac{\varepsilon_k^{\top}}{\sigma_k}\,\alpha_k\frac{\partial D_\theta(y_{k-1})}{\partial\theta}\Big],
\label{eq:sfterm}
\end{equation}
and replacing $V_k(y_k)$ by the Monte-Carlo return $r(y_K)$, or by an advantage $V_k(y_k)-b(y_{k-1})$ with any baseline $b$, leaves the expectation unchanged because $\mathbb E[\varepsilon_k\mid y_{k-1}]=0$. Summing~(\ref{eq:sfterm}) over $k$ and over the positions of $y_k$ is~(\ref{eq:sf}).

\paragraph{Step 2: the pathwise form of the same term.}
Condition on $y_{k-1}$ and use the reparameterisation $y_k=\alpha_kD_{\theta_k}(y_{k-1})+\sigma_k\varepsilon_k$. The objective depends on $\theta_k$ only through $y_k$, and the downstream value $V_k$ does not depend on $\theta_k$, so
\begin{equation}
\frac{\partial J}{\partial\theta_k}=\mathbb E\Big[\frac{\partial}{\partial\theta_k}V_k\big(\alpha_kD_{\theta_k}(y_{k-1})+\sigma_k\varepsilon_k\big)\Big]=\mathbb E\Big[\sum_{i=1}^{L}\frac{\partial V_k(y_k)}{\partial y_{k,i}}\,\alpha_k\,\frac{\partial D_\theta(y_{k-1})_i}{\partial\theta}\Big],
\label{eq:pwterm}
\end{equation}
which is exactly~(\ref{eq:pw}) with the critic $V_\phi$ replaced by the true value $V_k$. Exchanging derivative and expectation is legitimate for $k<K$ because, by~(\ref{eq:bellman}), $V_k$ is a Gaussian smoothing of a bounded function composed with $D_\theta$, hence as smooth as $D_\theta$ and with bounded gradient. For $k=K$ it is not: $V_K=r$ is piecewise constant in $y_K$ (the decode is an $\arg\max$ and the verifier a program), so its almost-everywhere gradient is zero and the pathwise identity carries no signal through the final transition. This is the precise sense in which ``the reward itself is not differentiable'' in \S\ref{sec:estimators}, and why the pathwise route needs a value estimate on \emph{intermediate} states rather than the reward on the final one. \emph{Conclusion of Steps 1--2:} with an exact critic, (\ref{eq:pw}) evaluated at step $k$ is an unbiased estimate of $\partial J/\partial\theta_k$, the $k$-th transition's contribution to the gradient of the expected reward.

\paragraph{Step 3: (\ref{eq:pw}) is~(\ref{eq:sf}) with the noise integrated out.}
The two forms of $\partial J/\partial\theta_k$ are related by Gaussian integration by parts (Stein's identity): for $\varepsilon\sim\mathcal N(0,I_n)$ and any $f$ that is almost-everywhere differentiable with $\mathbb E|f|,\mathbb E\|\nabla f\|<\infty$,
\begin{equation}
\mathbb E\big[f(\mu+\sigma\varepsilon)\,\varepsilon\big]=\sigma\,\mathbb E\big[\nabla f(\mu+\sigma\varepsilon)\big].
\label{eq:stein}
\end{equation}
With $f=V_k$, $\mu=\alpha_kD_\theta(y_{k-1})$ and $\sigma=\sigma_k$, (\ref{eq:stein}) turns the score-function factor $V_k(y_k)\,\varepsilon_k/\sigma_k$ of~(\ref{eq:sfterm}) into the pathwise factor $\nabla_yV_k(y_k)$ of~(\ref{eq:pwterm}), conditionally on $y_{k-1}$. The two estimators therefore have the same expectation and differ only in variance, and the difference is visible coordinate by coordinate: the score-function factor at coordinate $(k,i)$ is $V_k(y_k)\,\varepsilon_{k,i}/\sigma_k$, a random variable of variance $\approx\mathbb E[V_k^2]/\sigma_k^2$ whether or not the coordinate affects the reward, whereas the pathwise factor is $\partial V_k/\partial y_{k,i}$, which is identically zero on every coordinate the value does not depend on. This is the variance argument of Appendix~\ref{app:variance} in one identity: the noise factor that integration by parts removes is the same factor that makes~(\ref{eq:sf}) average each reward bit over the $K\cdot L\cdot d$ coordinates of the trajectory.

\paragraph{Step 4: one random flow step.}
\jpc{} draws one $k$ uniformly from a window $\mathcal W=\{k_0,\dots,k_1\}$ and applies~(\ref{eq:pw}) at that step. Taking the expectation over $k$,
\begin{equation}
\mathbb E_{k\sim\mathcal U(\mathcal W)}\big[\hat g_{\mathrm{PW},k}\big]=\frac1{|\mathcal W|}\sum_{k\in\mathcal W}\frac{\partial J}{\partial\theta_k},
\label{eq:window}
\end{equation}
so the one-step update is, up to the constant $|\mathcal W|$, an unbiased Monte-Carlo estimate of the window's share of~(\ref{eq:sumover}). The transitions outside the window contribute little to~(\ref{eq:sumover}): at the earliest steps the state is dominated by noise and $V_k$ is nearly flat (the pooled critic value is at chance for $k\lesssim K/8$, Table~\ref{tab:critic_quality}), and at the last steps the state is committed and $V_k\to r$ is piecewise constant, with zero gradient almost everywhere. Restricting to the window therefore drops terms that are small or exactly zero, and it is what makes the update's memory independent of $K$: the downstream dependence of the reward on $y_k$ is carried by $V_k$ in expectation, through~(\ref{eq:bellman}), not by back-propagating through the remaining $K-k$ network evaluations as in reward back-propagation for image diffusion. Differentiating through $m$ consecutive transitions instead would recover the sum of $m$ consecutive terms of~(\ref{eq:sumover}) by the chain rule through the intermediate $D_\theta$; we found no benefit from this (Appendix~\ref{app:results}), consistent with the terms being individually estimable.

\paragraph{Step 5: the learned critic, and why it must be trained jointly.}
Replacing $V_k$ by the critic $V_\phi$ in~(\ref{eq:pwterm}) adds the bias
\begin{equation}
\mathbb E\Big[\sum_i\frac{\partial\,(V_\phi-V_k)(y_k)}{\partial y_{k,i}}\,\alpha_k\frac{\partial D_\theta(y_{k-1})_i}{\partial\theta}\Big],
\label{eq:bias}
\end{equation}
so the update ascends $J$ to the extent that the critic's slope matches the true value's slope on the states the current policy visits. Two things move that target as $\theta$ changes: the distribution of $y_k$ under the policy, and $V_k$ itself, which by~(\ref{eq:bellman}) is a function of the downstream policy, so a critic fitted to an older policy is wrong even at the same state. Both are why the critic is refitted on the current policy's own rollouts after every actor step (\S\ref{sec:method}); training it by binary cross-entropy on the verifier outcome fits $\sigma(V_\phi)$ to $\mathbb E[r\mid y_k]$, i.e.\ to~(\ref{eq:valuedef}) up to the logistic link. Equation~(\ref{eq:bias}) also describes the failure mode: wherever $V_\phi-V_k$ has a positive slope the actor ascends $V_\phi$ while the true expected reward is flat or falls, which is the frozen-critic exploit of Appendix~\ref{app:frozen}, and the reason the recipe stops at the held-out peak.

\subsection{Token commitment along the trajectory}
\label{app:commit}
For each generated position we record the first step at which the denoiser's prediction $\hat x_k$ decodes to the final token and never changes afterwards. On Countdown, 86\% of positions commit at step 0 (median 0, 90th percentile 2), identically for $\gamma\in\{4,8,14\}$; the remaining contested positions commit uniformly along the trajectory. Word decisions therefore live in the first forward pass and in the self-conditioned prediction channel, not in the noise.
